% Nota de gerador_contorno. GERADA por flim_al/tabelas.py.
\footnotesize $\Delta$ é a variação de F$\beta$ na validação ao acrescentar \textbf{uma} anotação, do braço de contorno menos o braço de superpixel. Os dois sorteiam o candidato e gastam os mesmos ~300 px: a \textbf{única} diferença é o gerador da lista, e por isso a comparação isola o gerador e não diz nada sobre qual escore usar. A unidade é a \textbf{semente}; os dois decodificadores de uma semente compartilham o encoder e entram colapsados por média. Pior caso e dispersão foram declarados no pré-registro, com o mecanismo escrito antes de rodar: o candidato de contorno é pequeno e local, acrescenta menos filtros, e o ganho marginal mede correlação negativa entre filtros acrescentados e ganho. O nulo do BraTS é \textbf{previsto}: a arquitetura dele fixa o banco em 8 filtros por camada, os dois braços acrescentam exatamente 24, e o deslocamento que a intervenção evita não pode ocorrer. Semente cuja base degenerou (F$\beta$=0) sai, porque $\Delta$ a partir de zero não pode ser negativo. \textbf{O desfecho primário vem com três testes}: t pareado, teste de sinais e permutação exata de sinais. Com 3 e 5 sementes a normalidade das diferenças não é verificável, e o t sozinho não autorizaria conclusão; onde os três concordam, como no conjunto de parasitas, a conclusão não depende da suposição. \textbf{ATENÇÃO ao pior caso da conjuntivite}: o valor de 0,0480 é bruto e \textbf{não sobrevive} à análise de sensibilidade que casa o número de candidatos entre os braços, na qual passa a 0,1032. O gerador de contorno entregou UM único candidato naquela semente, e mínimo de um sorteio contra mínimo de oito não é comparação. Esse desfecho \textbf{não deve ser lido como significativo}.
\normalsize
